\documentclass[11pt]{article}
\usepackage[T1]{fontenc}
\usepackage[utf8]{inputenc}
\usepackage{lmodern,amsmath,amssymb,amsthm,mathtools}
\usepackage[margin=1in]{geometry}
\usepackage{microtype,booktabs,array,longtable,enumitem,xcolor,listings,xurl}
\usepackage[colorlinks=true,linkcolor=blue!45!black,urlcolor=blue!45!black,citecolor=blue!45!black]{hyperref}
\usepackage{fancyhdr}
\hypersetup{pdftitle={Concise iPRA: scope reuse and mathematical limits},pdfauthor={Research report prepared for Vladimir Reshetnikov}}
\pagestyle{fancy}\fancyhf{}\fancyhead[L]{Concise iPRA: scope reuse and mathematical limits}\fancyhead[R]{\thepage}
\setlength{\headheight}{14pt}
\setlength{\emergencystretch}{2em}
\setlist{itemsep=3pt,topsep=5pt}
\lstset{basicstyle=\ttfamily\footnotesize,breaklines=true,columns=fullflexible,keepspaces=true,frame=single,rulecolor=\color{black!20}}
\newtheorem{theorem}{Theorem}[section]
\newtheorem{lemma}[theorem]{Lemma}
\newtheorem{proposition}[theorem]{Proposition}
\newtheorem{corollary}[theorem]{Corollary}
\theoremstyle{definition}\newtheorem{definition}[theorem]{Definition}
\theoremstyle{remark}\newtheorem{remark}[theorem]{Remark}
\newcommand{\iPRA}{\mathrm{iPRA}}
\newcommand{\FV}{\operatorname{FV}}
\newcommand{\Teq}{T_{\mathrm{eq}}}
\newcommand{\Q}{\mathsf Q}
\newcommand{\D}{\mathsf D}
\newcommand{\C}{\mathcal C}
\newcommand{\llbracket}{\mathopen{[\![}}
\newcommand{\rrbracket}{\mathclose{]\!]}}
\title{Concise iPRA: a smaller validated checker\\[5pt]
\large Exact scope reuse, a product obstruction, and a partial uniqueness theorem}
\author{Research report prepared for Vladimir Reshetnikov}
\date{November 11, 2026}
\begin{document}
\maketitle
\begin{abstract}
This report gives a modest implementation improvement to the current
ConciseIPRA presentation. The full Eq/Q/D standalone checker falls from
897 to 881 normalized Python bytes, from 401 to 393 significant tokens,
and from 468 to 460 AST nodes. Its exact integer/list certificate contract,
public formula language, and accepted inputs are preserved. The Eq/D
public variant falls from 878 to 862 bytes. Gzip-9 sizes are reported only
for information. Both changes are proved to preserve acceptance and are
checked against the frozen incoming implementations.

The mathematical axiomatization and its existing bi-interpretation with
iPRA are unchanged. No new smaller axiom basis, global minimum, or complete
new Lean verification is claimed. In particular, these validated profiles
do not beat the separate 520-byte guarded profile with its different
input contract.

Two mathematical investigations are recorded precisely. A closed native
formula separates the existing case rule from a tempting direct-target
replacement, using a coordinatewise model on $\mathbb N^2$. A one-premise
iteration-uniqueness rule is proved to imply induction for zero equations
and, more generally, for predicates admitting a term-defined retraction.
The equality-reflection step needed to make that second route a complete
replacement for scalar equation induction remains open in this work.
\end{abstract}
\newpage
\tableofcontents
\newpage

\section{What is established, and what is not}

The requested mathematical goal is a simpler axiomatization bi-interpretable
with full first-order intuitionistic primitive recursive arithmetic.
The present work establishes a smaller complete \emph{implementation of the
existing presentation}, under the proof-checker-size criterion. It does
not establish a new smaller mathematical axiomatization. This distinction
is essential to interpreting the result.

The repository comparison is pinned to commit
\begin{center}\small\texttt{56249a4263cae9aad9ba87d23d0ee6d4ff6483bc}.\end{center}
The inspected directory is
\path{Logic/PrimitiveRecursiveArithmetic/ConciseIPRA}. Besides the main
README and selected Lean sources, the comparison includes all three ZIP
archives present in \path{docs/incoming} at this commit:
\texttt{ipra\_checked\_refinement}, \texttt{ipra\_one\_sided\_transport},
and \texttt{ipra\_validated\_core}, all dated 2026-11-11.
The relevant exact controls are copied into this package; their provenance
and source hashes are recorded in \path{PROVENANCE.json}.

This matters because the main README's compact envelope is not the newest
result in every interface. The incoming checked-refinement package already
has a 520-byte compact guarded checker, and the incoming validated-core
package already has the 897-byte fully typed Eq/Q/D checker used here.
One-sided crossed equality transport is also already present and is not
claimed as a new result. These facts were read from the actual snapshot,
not inferred from package names \cite{readme,validated,checked,onesided}.

\subsection{Matched measurements}

\begin{center}
\begin{tabular}{@{}lrrrr@{}}
\toprule
Complete program & Bytes & Tokens & AST nodes & Gzip-9\\
\midrule
Full Eq/Q/D: incoming & 897 & 401 & 468 & 505\\
Full Eq/Q/D: this package & \textbf{881} & \textbf{393} & \textbf{460} & 500\\
Eq/D: incoming & 878 & 390 & 454 & 498\\
Eq/D: this package & \textbf{862} & \textbf{382} & \textbf{446} & 493\\
\bottomrule
\end{tabular}
\end{center}

The full profile saves 16 bytes, approximately 1.78 percent. The Eq/D
profile saves 16 bytes, approximately 1.82 percent. Each saves eight
significant tokens and eight AST nodes. The same diagnostic text is kept;
none of the saving comes from shortening an error message.

The 520-byte profile has a guarded public language and assumes that its
advertised goal is canonical. The 788-byte Eq/D goal-checking profile
does not enforce the same exact whole-certificate type contract. Neither
is a matched comparator for the fully typed programs above. This report
makes no claim that 881 or 862 is a new minimum across those different
contracts. Reusing a kernel scope check already appeared in that weaker
Eq/D interface; the present result composes it with the stronger
whole-certificate contract.

\section{Complete inherited mathematical presentation}
\label{sec:presentation}

The source transformation leaves the proof calculus unchanged. For
clarity, its complete mathematical inventory is stated here. The
implementation uses the source-restricted form of the inherited
Choice-DI calculus. Its equivalence to the unrestricted native calculus
is an inherited result, not a new theorem of this report
\cite{readme,native,inherited}.

\subsection{Terms, formulas, and ordinary readings}

There is one numerical sort, with terms
\[
 t ::= x\mid 0\mid S(t)\mid I(n,a;u.h).
\]
Only the update body $h$ binds $u$. The counter $n$ and initializer $a$
are outside that binding. Arbitrary fixed numerical parameters may occur
in the body. Terms are finite syntax; variables do not range over
function descriptions.

Native formulas are
\[
 A ::= (s=t)\mid \Q_y A\mid \D_{x,y} A.
\]
The two bound variables of $\D$ are distinct. All substitutions are
capture-free, with bound variables renamed when necessary. The designated
contradiction is $0=S0$. The ordinary readings are
\begin{align*}
 (s=t)^\circ&:=(s=t),\\
 (\Q_yB)^\circ&:=B^\circ(0)\to\forall y\,B^\circ(y),\\
 (\D_{x,y}B)^\circ&:=\exists x\bigl(B^\circ(x,0)\to
                                  \forall y\,B^\circ(x,y)\bigr).
\end{align*}
These readings apply also to atomic matrices; a progressive proof
instruction does not change a formula's meaning.

\subsection{Every primitive axiom and rule}

Contexts $\Gamma$ are finite reusable lists of native formulas. The
computation axioms, in any retained context, are
\begin{align}
 I(0,a;u.h)&=a,\tag{I0}\label{eq:I0}\\
 I(Sn,a;u.h)&=h[I(n,a;u.h)/u].\tag{IS}\label{eq:IS}
\end{align}
The other rules are exactly the following.

\paragraph{Hypothesis.}
If $A$ is retained in $\Gamma$, then $\Gamma\vdash A$.

\paragraph{Equality transport.}
For every native formula template $F(v)$,
\[
 \frac{\Gamma\vdash s=t\qquad\Gamma\vdash F(s)}
      {\Gamma\vdash F(t)}.
 \tag{Tr}
\]
The source-restricted implementation requires that $s$ is not the
\emph{literal syntax tree} $0$. It may denote zero: for example,
$I(0,0;u.0)$ is an allowed source. All omitted literal-zero instances
are derivable in this calculus. The inherited Lean kernel instead uses
unrestricted transport for equation templates only and derives general
formula transport; the resulting sequent relations coincide.

\paragraph{Closure elimination.}
\[
 \frac{\Gamma\vdash\Q_yB(y)\qquad\Gamma\vdash B(0)}
      {\Gamma\vdash B(t)}.
 \tag{QE}
\]

\paragraph{Paired-binder elimination.}
\[
 \frac{\Gamma\vdash\D_{x,y}B(x,y)\qquad
       \Gamma,\Q_yB(b,y)\vdash A}
      {\Gamma\vdash A}.
 \tag{DE}
\]
Here $b$ is fresh for $\Gamma$, $A$, and every fixed parameter of $B$;
the displayed local assumption is discharged.

\paragraph{Introduction by cases.}
\[
 \frac{\Gamma,B(t,0)\vdash B(t,Sz)}
      {\Gamma\vdash\D_{x,y}B(x,y)}.
 \tag{DI-case}
\]
Here $z$ is fresh for $\Gamma$, $t$, and the fixed parameters of $B$.
The matrix $B$ may be any native formula.

\paragraph{Introduction by equation progress.}
\[
 \frac{\Gamma,E(t,z)\vdash E(t,Sz)}
      {\Gamma\vdash\D_{x,y}E(x,y)}.
 \tag{DI-progress}
\]
The same freshness condition applies, and $E(x,y)$ must be a literal
equation, not an equation inside a binder. The local assumption is
discharged. There are no other primitive logical, induction, successor,
or equality axioms. In particular, reflexivity and explosion are derived.

The Lean presentation with direct atomic closure replaces the last rule
by $\Gamma,E(z)\vdash E(Sz)\Rightarrow\Gamma\vdash\Q_yE(y)$.
The unused-witness construction proves the equivalence of these two
forms. The present Python transformation changes neither form nor any
certificate instruction.

\subsection{Derived induction and numerical cases}

For any literal equation $E(y)$, progress introduction with an unused
witness coordinate, followed by DE and QE, gives
\[
 \frac{\Gamma\vdash E(0)\qquad\Gamma,E(z)\vdash E(Sz)}
      {\Gamma\vdash E(t)}.
 \tag{EI}
\]
For an arbitrary native formula $B(y)$, the same unused-witness
construction using DI-case gives
\[
 \frac{\Gamma\vdash B(0)\qquad\Gamma\vdash B(Sz)}
      {\Gamma\vdash B(t)}.
 \tag{Cases}
\]
The second premise may ignore the retained base assumption. Freshness is
as above. These two principles have different mathematical roles:
Section~\ref{sec:product-obstruction} explains why EI does not replace Cases.

As a simple equality example, put $r=I(0,t;u.0)$. I0 gives $r=t$,
and transport with template $v=t$, using that equation as both premises,
gives $t=t$. The source is iterator-headed, so the source restriction is
respected.

\subsection{All ordinary logical operations are explicit macros}

Let $C(v,a,b)=I(v,a;u.b)$, with $u$ fresh. To define a formula selector
without adding a connective, first define $L(v,s,t,B)$ by recursion on $B$:
\begin{align*}
 L(v,s,t,r=q)&:=\bigl(C(v,s,r)=C(v,t,q)\bigr),\\
 L(v,s,t,\Q_yB)&:=\Q_y L(v,C(y,0,s),C(y,0,t),B),\\
 L(v,s,t,\D_{x,y}B)&:=\D_{x,y}L(v,C(y,0,s),C(y,0,t),B).
\end{align*}
All newly bound variables are fresh for external arguments. Define
$G(v,A)=L(v,0,0,A)$. Next define $\operatorname{Sel}$ by recursion on its
first formula argument:
\begin{align*}
 \operatorname{Sel}(v,s=t,B)&:=L(v,s,t,B),\\
 \operatorname{Sel}(v,\Q_yA,B)&:=\Q_y\operatorname{Sel}(v,A,G(y,B)),\\
 \operatorname{Sel}(v,\D_{x,y}A,B)&:=\D_{x,y}\operatorname{Sel}(v,A,G(y,B)).
\end{align*}
The inherited selector lemmas prove that $G(0,A)$ is a theorem,
$G(Sv,A)$ is equivalent to $A$, and $\operatorname{Sel}(0,A,B)$ and
$\operatorname{Sel}(Sv,A,B)$ are equivalent to $A$ and $B$, respectively.

With $\operatorname{par}(y)=I(y,0;u.C(u,1,0))$, the logical macros are
\begin{align*}
 \bot&:=0=S0,\\
 A\to B&:=\Q_y\operatorname{Sel}(y,A,B),\\
 A\land B&:=\Q_yG(y,\operatorname{Sel}(\operatorname{par}(y),A,B)),\\
 \exists x A&:=\D_{x,y}G(y,A),\\
 A\lor B&:=\exists v\,\operatorname{Sel}(v,A,B),\\
 \forall x A&:=A(0)\land\Q_xA.
\end{align*}
These are recursive syntactic abbreviations. The existing proofs establish
the full intuitionistic natural-deduction rules for them, including
assumption discharge, the two quantifier eigenvariable conditions, and
explosion. There is no appeal to classical excluded middle for arbitrary
formulas. The complete inherited choice-basis proof is supplied in
\path{reference/mathematics/inherited_choice_basis.tex}.

\section{Why the iPRA bi-interpretation is preserved}
\label{sec:bridge}

It is useful to separate a theorem about mathematical theories from a
theorem about program acceptance. The current source already establishes
the former; the new proof in Section~\ref{sec:preservation} establishes
the latter for these modifications.

\subsection{The intermediate ordinary theory}

Let $\Teq$ use ordinary intuitionistic first-order logic with actual
equality and the same scalar terms. Its nonlogical axioms and schema are
I0, IS, numerical coverage
\[
 \forall x\,(x=0\lor\exists y\,x=Sy),
\]
separation $0=S0\to\bot$, and contextual induction EI for literal
equations. The ordinary interpretations of $\Q$ and $\D$ send every
native rule to a valid derivation in $\Teq$. Conversely, the macros above
translate the ordinary logical rules and the nonlogical principles into
native derivations. Both formula translations are faithful and their
composites are provably equivalent to the original formulas.

The scalar-to-iPRA interpretation compiles each fixed scalar term into
a fixed primitive-recursive description. The iterator is a primitive
recursion whose update ignores the recursion counter. This is structural
compilation of a finite expression, not a postulated universal
primitive-recursive evaluator.

For the reverse interpretation, the inherited scalar development first
constructs a predecessor satisfying $P(0)=0$ and $P(Sn)=n$. It then
constructs a faithful equality characteristic, proves quantifier-free
induction by Boolean reflection, constructs pairing and projections, and
compiles arbitrary fixed primitive-recursive descriptions into scalar
terms. The order prevents circular reliance on pairing to obtain
predecessor or on equality reflection to justify its own induction.
The supplied source files \path{scalar_predecessor.tex},
\path{equality_reflection.tex}, and \path{primitive_recursion.tex} contain
these inherited written arguments.

For each fixed primitive-recursive description, uniqueness of its
recursion equations proves that translating its scalar compilation back
to iPRA yields the original function. Structural induction on scalar terms,
including uniqueness of each fixed iterator recurrence, proves the other
term round trip. The formula statements then follow by structural induction
and the native macro equivalences.

\subsection{Exact strength of the inherited result}

Write $J$ for native-to-iPRA translation and $K$ for the reverse.
The source proves, for arbitrary finite contexts and open formulas,
\begin{align*}
 \Gamma\vdash_{\C}A&\quad\Longleftrightarrow\quad
       J\Gamma\vdash_{\iPRA}J(A),\\
 \Delta\vdash_{\iPRA}\varphi&\quad\Longleftrightarrow\quad
       K\Delta\vdash_{\C}K(\varphi).
\end{align*}
It also proves $A$ and $KJ(A)$ interderivable and
$\iPRA\vdash\varphi\leftrightarrow JK(\varphi)$, with parameters.
Both numerical domains remain whole and both numerical composites have
identity graph $x=y$. This is bi-interpretability, with the stronger
identity-domain property, rather than merely equiconsistency or
agreement about standard numerical calculations.

The named Lean statements include \texttt{native\_iff\_pra},
\texttt{pra\_iff\_native}, \texttt{native\_roundTrip},
\texttt{pra\_native\_roundTrip}, and
\texttt{native\_inconsistent\_iff} \cite{bridge}. The comparison with
the reports' rule inventory is \texttt{iprf\_iff\_nprf}
\cite{inherited}. Selected exact source snapshots are included for review.
They are not a standalone copy of the entire Lean project. This work
inspected those sources; it did not rebuild the inherited Lean chain or
formalize the Python-to-Lean correspondence.

In particular, the unchanged calculus proves $0=S0$ exactly when iPRA
proves $\bot$. The new programs accept exactly the same input strings as
their matched old programs in the abstract execution sense stated below.
Thus they inherit exactly the same represented theorem set and the same
bi-interpretation. This last step relies on the inherited checker
correctness argument; acceptance preservation does not independently
re-prove that whole argument.

\section{The two implementation changes and their proofs}
\label{sec:preservation}

\subsection{Input contract and trusted boundary}

Each selected checker is one complete Python module. It imports only the
standard JSON library, reads standard input, and accepts by normal exit
with output exactly \texttt{ACCEPT} followed by one newline. Any exception
or other result is rejection. No axiom database, certificate producer,
measurement utility, or theorem library is imported by the checker.
The code does not use Python assertions for proof obligations, so
\texttt{python3 -O} retains those obligations.

The input is a nonempty JSON list. The final element advertises the
conclusion. Earlier elements supply four-field packets consumed from
right to left. An unused prefix is allowed but is still fully checked to
contain only exact integers and lists. Booleans, floats, strings,
null, and dictionaries anywhere in that prefix are rejected. The full
profile permits canonical Eq/Q/D conclusions; the smaller profile
permits Eq/D conclusions. Internal closures remain available in both.
The smaller public language still expresses the whole theory: replace
$\Q_yB$ by $\D_{x,y}B$ with $x$ absent from $B$, recursively throughout
the formula. On the nonempty numerical domain their ordinary readings
agree, and DE together with DI-case and QE proves the native
interderivability. Thus the Eq/D profile does not discard the meaning
of closure formulas.

\subsection{Sorted arity syntax}

At a formula position, arities one, two, and three mean $\D$, equality,
and $\Q$, respectively. A public $\Q$ is represented by three identical
copies of its matrix. At a numerical position, arities one and three
mean successor and iterator. Canonical zero is the integer zero.
The interpretation is sorted: a ternary formula and a ternary numerical
term have different meanings.

The inherited encoding gives every protected numerical slot width two.
At a numerical position under $p$ protected slots, the scope depth is
$e=3+2p$. Canonical variable coordinates are the odd numbers
$3,5,\ldots,2p+1$. Under successor, two dummy zero slots are introduced;
under an iterator, one state slot is introduced in its body and a dummy
zero slot in its counter and initializer. These dummy slots are part of
the published decoder, not extra free variables. Under $\D$ two actual
numerical slots are introduced, and under $\Q$ one is introduced.

Every list edge updates the opener depth by
\[
                         e\longmapsto e+5-\operatorname{len}(a).
\]
The public validator uses a depth $d$ larger by two, starting with
$d=2$ at the formula root. Thus it follows the same increments and
$d=e+2$ at every node. The complete original sorted decoder and the
handling of administrative auxiliary arrays are included in
\path{reference/mathematics/inherited_arity_checker.tex}. No new meaning
for auxiliary arrays is introduced by this report.

\subsection{Reuse the scope check already performed by the kernel}

At a numerical leaf, the old public validator permits zero or an odd
integer satisfying
\[
                          2<a<d-2.
\]
The new validator keeps exact integer typing, numerical position,
oddness, and the lower bound, but replaces that chained comparison by
$a>2$. The omitted upper bound is enforced by the first unchanged
operation of the proof kernel.

\begin{lemma}[Initial opening enforces the omitted bound]
\label{lem:scope}
For a public goal that passes the new shape-and-type validator, its
initial no-image opening succeeds if and only if every numerical slot
satisfies the old upper bound. When it succeeds, the opening returns the
same raw goal.
\end{lemma}
\begin{proof}
At the first kernel call the freshness floor is $k=0$. The opener
traverses the complete goal with a list comprehension before it checks
context membership or reads an instruction. At a scalar it retains the
value precisely when $0\leq a<e$. On canonical numerical positions,
zero satisfies this condition. For positive allowed slots, failure is
exactly $a\geq e=d-2$, the omitted bound.

At such a failure, numerical parity makes the opener attempt to insert
its default image. That default is the checking function itself, used
as an invalid-image sentinel. Evaluating the opener's type-dispatch
operation on a function raises \texttt{TypeError}. The checker therefore
rejects before any proof packet can be consumed. No hypothesis can
hide the failure, because the complete initial traversal is evaluated
before membership is tested.

Conversely, when all slots satisfy the old bounds, the opener retains
every leaf and recursively reconstructs exactly the same lists. Structural
induction on the canonical goal proves that its output equals its input.
\end{proof}

\begin{theorem}[Scope reuse preserves acceptance]
Replacing only $2<a<d-2$ by $a>2$ in either matched public validator
preserves the successful input set.
\end{theorem}
\begin{proof}
The raw prefix validator is unchanged, as are every arity, type, parity,
lower-bound, and public-language check. A goal passing the old validator
passes the new one and produces the identical first kernel state. A goal
newly admitted by the relaxed validator fails the initial opening by
Lemma~\ref{lem:scope}. All later kernel operations are byte-for-byte
unchanged. These observations cover both directions.
\end{proof}

The success-set statement abstracts from physical memory and recursion
limits. The location and text of a rejection need not agree. The theorem
asserts that a malformed scope cannot become a successfully checked proof,
not that both implementations allocate the same amount of memory.

\subsection{Split the certificate directly}

The second change replaces
\begin{lstlisting}[language=Python]
q = json.load(open(0))
a = q.pop()
\end{lstlisting}
by
\begin{lstlisting}[language=Python]
*q, a = json.load(open(0))
\end{lstlisting}

\begin{lemma}[Envelope preservation]
The direct split preserves the successful input set under the retained
validators.
\end{lemma}
\begin{proof}
For every nonempty JSON list, it produces precisely the same prefix
and last element as the old pair of statements. The checker never tests
the identity of the certificate's list containers. Empty iterables and
noniterable JSON roots reject.
The remaining alternative JSON roots are nonempty strings and objects:
unpacking either supplies a string as the alleged goal (an object is
iterated through its string keys). A string cannot pass the unchanged
goal-type validator. Thus no formerly rejecting root becomes accepted.
\end{proof}

\begin{corollary}[Exact matched-program preservation]
The delivered full checker and Eq/D checker have the same successful
input sets as their respective frozen incoming controls, subject to the
abstract resource qualification above. In particular, both retain exact
whole-certificate integer/list validation, including unused prefixes.
\end{corollary}

\noindent\begin{minipage}{\textwidth}
\subsection{Complete primary program}

The following is the entire 881-byte normalized full-language checker.
It is printed as an audit aid; the executable file is authoritative.
\lstinputlisting[language=Python,basicstyle=\ttfamily\scriptsize,showstringspaces=false]{code/checker_full.py}
The readable standalone implementation, \path{code/checker_readable.py},
uses descriptive names and comments. It is supplied for review and was
included in the differential process checks. Its larger source size is
not substituted for the counted program.
\end{minipage}

\section{Reproduction, measurements, and validation}

From the extracted package root, run
\begin{lstlisting}
python3 code/checker_full.py < examples/computation.json
python3 -O code/checker_full.py < examples/computation.json
python3 reproduce.py
python3 build_report.py
\end{lstlisting}
The first two commands print \texttt{ACCEPT}. The reproduction command is
offline and uses the Python standard library. The document build requires
\texttt{pdflatex} and the standard LaTeX packages listed in this source.

The frozen measurement runtime was CPython 3.12.14. Normalize a whole
module by \texttt{ast.unparse(ast.parse(source))} plus a final linefeed.
Bytes mean the UTF-8 size of that normalized source. Significant tokens
exclude comments, encoding/end markers, indentation, dedentation, and
newline tokens. AST nodes count every node yielded by
\texttt{ast.walk}. Gzip is level nine with \texttt{mtime=0}, including
the gzip wrapper. Its size is informational, not an optimization target.
The normalized gzip artifacts and exact source SHA-256 values are saved
in \path{evidence/}.

The packaged verification was actually executed. It recorded 13,019
scope-preflight comparisons \emph{per interface}, with 3,096 valid
full-language examples and 2,233 valid Eq/D examples, and no discrepancy.
It also ran 606 complete checker processes across the old and new
programs, the readable full implementation, and ordinary and optimized
Python, with no unexpected outcome. Thirty additional targeted controls
cover non-list JSON roots, high positive slots, negative scope escapes,
Boolean/float/null aliases, malformed unused prefixes, and valid inert
integer/list prefixes.

The preflight checks exercise the actual first kernel opening with a
sentinel instruction provider; reaching the proof boundary is
distinguished from accepting a proof. The process checks exercise the
complete input-to-output programs. The detailed records are in
\path{evidence/validation.json}; the structural proofs above supply the
general preservation argument that finite testing cannot supply.

\section{Why equation induction cannot replace numerical cases}
\label{sec:product-obstruction}

A tempting simplification changes the case introduction from a proof of
$B(t,Sz)$ under $B(t,0)$ to a proof of $B(t,z)$ under $B(t,0)$, with $z$
fresh.  The modified rule is sound for the ordinary conditional-closure
reading without any numerical coverage axiom.  It is nevertheless strictly
weaker when combined with the remaining native rules and literal-equation
induction.  This section gives both a semantic obstruction and a closed
native formula separating the two calculi.  These results do not assert
that every possible alternative interpretation of the weaker theory is
impossible.

\subsection{A product model of all the equational arithmetic}

Let $\mathcal A$ have numerical domain $\mathbb N\times\mathbb N$.  Interpret
zero, successor, and iteration coordinatewise:
\[
 0^{\mathcal A}=(0,0),\qquad
 S^{\mathcal A}(a,b)=(a+1,b+1).
\]
For each finite iterator body, the first output coordinate is the usual
natural-number iteration of its first-coordinate interpretation, and the
second coordinate is defined in the same way.  Equality is actual equality
of ordered pairs.  Use ordinary two-valued first-order semantics; a
classical structure is also a sound test of intuitionistic derivability.

More explicitly, if $\rho$ assigns pairs to the free variables, write
$\rho_i$ for its $i$th-coordinate assignment.  Structural induction on a
scalar term $t$ gives
\begin{equation}
 \llbracket t\rrbracket_{\mathcal A,\rho}
 =\bigl(\llbracket t\rrbracket_{\mathbb N,\rho_1},
         \llbracket t\rrbracket_{\mathbb N,\rho_2}\bigr).
 \label{eq:product-term-factorization}
\end{equation}
The iterator case follows by the defining coordinatewise interpretation,
including the coordinatewise assignment to its bound state variable.
Consequently, for every literal equation $E(n,\vec p)$ and every fixed
parameter assignment, there are unary predicates $E_1,E_2$ on the standard
natural numbers such that
\[
 \mathcal A\models E((a,b),\vec p)
 \quad\Longleftrightarrow\quad E_1(a)\ \text{and}\ E_2(b).
\]
This assertion concerns equations.  Arbitrary formulas containing
implication or quantifiers need not factor this way.

Both iterator computation equations hold in $\mathcal A$, successor is
injective, and $0\ne S0$.  Moreover, every literal-equation induction
instance is valid.  To prove the latter claim, suppose its base and step
hold.  The base supplies $E_1(0)$ and $E_2(0)$.  Given $E_1(a)$, evaluate
the product step at $(a,0)$, using $E_2(0)$, to obtain $E_1(a+1)$.
Similarly, evaluating the step at $(0,b)$ gives
$E_2(b)\Rightarrow E_2(b+1)$.  Ordinary induction in each standard
coordinate yields $E_1(a)$ and $E_2(b)$ for all $a,b$.  This proves the
product induction conclusion.

The argument also proves soundness of contextual equation induction.
Fix an assignment making every retained assumption in $\Gamma$ true.
Freshness of the induction variable means that changing its pair value
does not change the truth of $\Gamma$.  The preceding argument therefore
applies to the base and step under that same assignment.  No restriction to
equational side assumptions is required.

Numerical coverage fails.  The element $(0,1)$ is neither $(0,0)$ nor a
successor: every successor has two positive coordinates.  Thus
\begin{equation}
 \mathcal A\not\models
 \forall x\bigl(x=0\lor\exists y\,x=Sy\bigr).
 \label{eq:product-no-coverage}
\end{equation}
Adding all valid predecessor computation equations does not remove this
obstruction: coordinatewise predecessor already satisfies
$P(0)=0$ and $P(Sx)=x$ in this model.

\subsection{The failure is visible inside the native language}

Interpret the native binders by their ordinary readings
\[
 Q_yB:=B(0)\to\forall y\,B(y),\qquad
 D_{x,y}B:=\exists x\bigl(B(x,0)\to\forall y\,B(x,y)\bigr).
\]
Hypothesis, equality transport, QE, and DE are sound in $\mathcal A$.
The two computations and equation progress are sound by the preceding
argument.  The proposed direct-target DI rule is sound as well: its fresh
variable already ranges over the whole domain.  Hence every theorem of
that modified calculus is valid in $\mathcal A$.

Write $1=S0$ and $C(v,a,b)=I(v,a;u.b)$, with $u$ absent from $b$.
Define the following native formulas, where the displayed $E$ is a single
literal equation:
\begin{align*}
 E(n,y,p)&:\quad
 C\bigl(p,C(n,y,1),C(n,1,0)\bigr)=0,\\
 P(n,p)&:=Q_yE(n,y,p),\\
 G(p)&:=Q_nP(n,p),\\
 F&:=Q_pG(p).
\end{align*}
The last formula is closed.

The original calculus proves $F$ without equation induction.  At a
successor $n=Sz$, computation makes $E(Sz,y,p)$ equivalent to
$C(p,1,0)=0$, independently of $y$.  Thus $P(Sz,p)$ follows by
conditional-closure introduction from that same equation as a retained
assumption.  The original case introduction, with an unused witness,
therefore proves $Q_nP(n,p)$: its successor premise is already proved and
its retained zero premise can be ignored.  A further such introduction,
now in the parameter $p$, proves $F$.  Only case DI, DE, QE, computations,
and equality reasoning are needed for these derived closure introductions.

In $\mathcal A$, take $p_*=(0,1)$.  Direct coordinatewise computation gives
\[
 \mathcal A\models E(n,y,p_*)
 \quad\Longleftrightarrow\quad
 n_1=0,\quad n_2>0,\quad y_1=0.
\]
Consequently $P(0,p_*)$ is true, because its antecedent is false.
Likewise $P(Sz,p_*)$ is true for every $z$.  But $P((0,1),p_*)$ is
false: its antecedent holds at $y=0$, whereas its universally quantified
consequent fails at $y=(1,0)$.  Thus $G(p_*)$ is false.

Finally, $G(0)$ is true.  At $p=0$ and $n=0$, the equation $E$ becomes
$y=0$, so $P(0,0)=Q_y(y=0)$ is false in this nontrivial structure.  That
false antecedent makes $G(0)$ true.  Since $G(p_*)$ is false, the closed
formula $F=Q_pG$ is false.  Soundness of the modified calculus in
$\mathcal A$ therefore proves that it cannot derive $F$, while the
original calculus does derive $F$.

\subsection{What remains possible}

The obstruction distinguishes algebraic Boolean identities from global
logical alternatives.  For example, equation induction proves
$t(C(n,a,b))=C(n,t(a),t(b))$ for each fixed term context $t$.
It can also prove that the Boolean selector value $b=C(n,0,1)$ satisfies
$C(b,0,1)=b$.  Such identities remain valid for the mixed Boolean pairs
$(0,1)$ and $(1,0)$.  They do not entail $b=0\lor b=1$.
Accordingly, improving the equational predecessor construction cannot by
itself recover the missing numerical case rule.

There is a separate, unfinished route through recursion uniqueness.
Much of the scalar predecessor bootstrap can be recast as uniqueness of
an explicitly displayed recurrence. Write
$Z(u)=C(u,1,0)$, $N(n)=I(n,0;u.Z(u))$, and
$a+b=I(b,a;u.Su)$, with $2a=a+a$. For example, if
$j(u)=Z(u)+2N(u)$, $M(n)=I(n,0;j)$, and
$G(n)=I(n,0;u.u+M(u)+1)$, the difficult comparison
$M(G(n))=N(n)$ can compare two recurrences with the common update
$k(v)=I(Sv,v;j)$.  Finite Boolean computations give
$k(N(n))=Z(N(n))$.  This is useful progress on that route, but it does not
complete equality reflection or establish equivalence of a new
uniqueness-only theory with iPRA; no such conclusion is claimed here.

\section{A conditional route from iterator uniqueness to induction}
\label{sec:uniqueness-research}

This section records an arithmetic research direction, not a completed
replacement for the delivered checker.  Its positive result is a precise
retraction argument: scalar iterator uniqueness implies induction for
every predicate for which a suitable numerical conditional is already
available.  In particular, it implies induction for equations of the
form $t(n)=0$ using only numerical coverage and the ordinary zero test.
The missing step is an independent equality-reflection bootstrap for
the scalar language.  No new bi-interpretation or smaller complete
checker is claimed for the proposed uniqueness rule.

\subsection{The rule and its elementary consequences}

Uniqueness of recursion as an alternative arithmetic principle is
classical: Tait states the equivalence for a richer language with
counter-aware primitive recursion and product types \cite{tait}.
That result alone does not settle the restricted scalar update
$h(u)$ considered here. No novelty is claimed for the general idea
of using recursion uniqueness.

Work, for this arithmetic analysis, in ordinary intuitionistic
first-order logic with actual equality.  The numerical terms are
\[
 t ::= x\mid 0\mid S(t)\mid I(n,a;u.h),
\]
with capture-free substitution; $u$ is bound only in $h$.  Retain the
two computation equations
\begin{align*}
 I(0,a;u.h)&=a,\\
 I(Sn,a;u.h)&=h[I(n,a;u.h)/u].
\end{align*}
Terms such as $f(z)$ and $h(u)$ below are fixed finite term templates,
not variables of a function sort.  Passive parameters are suppressed.

Consider replacing literal-equation induction by the rule
\[
 \frac{\Gamma\vdash f(Sz)=h(f(z))}
      {\Gamma\vdash f(t)=I(t,f(0);u.h(u))}.
 \tag{UR}\label{rule:scalar-ur}
\]
Here $z$ is fresh for $\Gamma$, the target term $t$, and all passive
parameters; the update template $h(u)$ does not contain $z$.  Its only
active argument is the previous state $u$.  The displayed premise is
proved in $\Gamma$ itself: there is no discharged induction hypothesis.
The seed $f(0)$ is determined by $f$, so no separate base premise or
seed certificate occurs in this rule.

\paragraph{Equation induction proves UR.}
Assume the premise of UR and induct on the literal equation
\[
 f(n)=I(n,f(0);u.h(u)).
\]
The base is the zero computation equation and symmetry.  At a successor,
the recurrence premise, the current induction equation, congruence,
and the successor computation equation give
\begin{align*}
 f(Sn)&=h(f(n))\\
      &=h(I(n,f(0);u.h(u)))\\
      &=I(Sn,f(0);u.h(u)).
\end{align*}
Thus UR is sound in the old arithmetic.  This direction needs neither
coverage nor a numerical equality test.

\paragraph{Uniqueness with a common initial value.}
UR implies the following useful derived rule.  If, with fresh $z$,
\[
 \Gamma\vdash f(0)=g(0),\qquad
 \Gamma\vdash f(Sz)=h(f(z)),\qquad
 \Gamma\vdash g(Sz)=h(g(z)),
\]
then $\Gamma\vdash f(t)=g(t)$.  Apply UR to $f$ and $g$ separately
and use equality congruence in the seed of the iterator.  The argument
is contextual: arbitrary retained first-order assumptions may occur
in $\Gamma$, provided they do not contain the eigenvariable.

\paragraph{The successor iterator is the identity.}
Taking $f(z)=z$ and $h(u)=Su$, the premise of UR is reflexivity.
Consequently
\[
 \Gamma\vdash I(t,0;u.Su)=t.
 \tag{Id}\label{eq:ur-identity}
\]
This identity is a consequence of UR, not an additional axiom.  In
particular, the retraction argument below does not silently assume
that a successor orbit already exhausts the numerical domain.

\subsection{Retractions closed under successor are the identity}

\paragraph{Retraction lemma.}
Let $e(x)$ be a fixed scalar term template.  Suppose that $\Gamma$
proves the following equations, with $x$ fresh:
\begin{align}
 e(0)&=0,\label{eq:ur-e0}\\
 e(e(x))&=e(x),\label{eq:ur-idempotence}\\
 e(S(e(x)))&=S(e(x)).\label{eq:ur-closed}
\end{align}
Then the computation equations, equality reasoning, and UR prove
$\Gamma\vdash e(t)=t$ for every target term $t$.

\paragraph{Proof.}
Define two finite term abbreviations
\[
 h(u)=e(S(e(u))),\qquad r(n)=I(n,0;u.h(u)).
\]
No new operation or function quantifier is introduced.  Idempotence
implies the two global commutation equations
\begin{align}
 e(h(u))&=h(u),\label{eq:ur-eh}\\
 h(e(u))&=h(u).\label{eq:ur-he}
\end{align}
For the first, substitute $S(e(u))$ for $x$ in
\eqref{eq:ur-idempotence}.  For the second, unfold $h$ and replace
$e(e(u))$ by $e(u)$ under $S$ and $e$.  The successor-closure
assumption is exactly
\[
 h(u)=S(e(u)).
 \tag{H}\label{eq:ur-h-successor}
\]

Now consider $f(n)=e(r(n))$.  Its base is
$f(0)=e(0)=0$.  Its successor equation is, using the iterator
computation and \eqref{eq:ur-eh}--\eqref{eq:ur-he},
\[
 f(Sn)=e(h(r(n)))=h(r(n))=h(e(r(n)))=h(f(n)).
\]
UR, followed by replacement of the equal seeds, therefore yields
\[
 e(r(n))=I(n,0;u.h(u))=r(n).
 \tag{R1}\label{eq:ur-er}
\]
This is the first substantive use of uniqueness; it establishes that
the entire $h$ orbit lies in the image of the retraction.

Equation \eqref{eq:ur-h-successor} and \eqref{eq:ur-er} now give
\[
 r(Sn)=h(r(n))=S(e(r(n)))=S(r(n)),\qquad r(0)=0.
\]
Apply UR again, this time to $r$ with update $u\mapsto Su$.
By \eqref{eq:ur-identity},
\[
 r(n)=I(n,0;u.Su)=n.
 \tag{R2}\label{eq:ur-r-identity}
\]
Combining \eqref{eq:ur-er} and \eqref{eq:ur-r-identity} gives
$e(n)=n$.  Every displayed calculation is an equality derivation
under the fixed context.  No induction is used in the proof.
\hfill$\square$

\subsection{Induction from term-definable conditionals}

The preceding lemma gives a useful sufficient condition for induction.
Let $E(n)$ be a fixed formula.  Suppose that the retained context
proves its base and that its step is derivable:
\[
 \Gamma\vdash E(0),\qquad
 \Gamma,E(z)\vdash E(Sz),
 \tag{P}\label{eq:ur-ind-premises}
\]
where $z$ is fresh.  Suppose also that $E$ is decidable in this
context and has a term-defined split: there is a fixed term $e(n)$
such that, uniformly in fresh $n$,
\begin{align}
 \Gamma&\vdash E(n)\lor\neg E(n),\label{eq:ur-dec}\\
 \Gamma,E(n)&\vdash e(n)=n,\label{eq:ur-true-branch}\\
 \Gamma,\neg E(n)&\vdash e(n)=0.\label{eq:ur-false-branch}
\end{align}
Then UR proves the induction conclusion $\Gamma\vdash E(t)$.

Here is the complete reduction.  A case split using
\eqref{eq:ur-dec} proves $E(e(n))$: in the $E(n)$ branch use
\eqref{eq:ur-true-branch}; in the other branch use
\eqref{eq:ur-false-branch} and the base $E(0)$.  The true-branch
identity, instantiated at $e(n)$, gives
\[
 e(e(n))=e(n).
\]
The step from \eqref{eq:ur-ind-premises}, instantiated at $e(n)$,
gives $E(S(e(n)))$.  Another use of the true-branch identity yields
\[
 e(S(e(n)))=S(e(n)).
\]
Finally, the base and the true-branch identity at zero yield
$e(0)=0$.  Thus all three premises of the retraction lemma hold.
It follows that $e(n)=n$, and transporting the already proved
$E(e(n))$ along this equality gives $E(n)$, then $E(t)$.

This result applies only to predicates whose displayed numerical
split has been justified.  It does not supply induction for arbitrary
quantified formulas.  In particular, no general principle of deciding
first-order formulas or converting them to numerical tests is being
assumed.

\subsection{An unconditional application to zero equations}

For this subsection also assume numerical coverage and separation:
\[
 \forall x\,(x=0\lor\exists y\,x=Sy),\qquad 0=S0\to\bot.
\]
Define the selector, as in the inherited scalar language, by
\[
 C(v,a,b)=I(v,a;u.b).
\]
Its two endpoint equations are immediate computation:
\[
 C(0,a,b)=a,\qquad C(Sv,a,b)=b.
\]
Coverage supplies the corresponding case distinction for every
argument.  Moreover, $Sx=0$ is contradictory: applying the term
$v\mapsto C(v,0,S0)$ gives $S0=0$, and separation then gives
$\bot$.  Hence every zero equation $q(n)=0$ is decidable, without
constructing a test for equality of two arbitrary terms.

Take $E(n)$ to be the literal equation $q(n)=0$ and set
\[
 e(n)=C(q(n),n,0).
\]
If $E(n)$ holds, the zero selector equation gives $e(n)=n$.
If $\neg E(n)$ holds, coverage for $q(n)$ eliminates its zero
alternative; in its successor alternative the other selector
equation gives $e(n)=0$.  Thus
\eqref{eq:ur-dec}--\eqref{eq:ur-false-branch} have direct proofs.
The conditional induction result proves the following derived rule:
\[
 \frac{\Gamma\vdash q(0)=0\qquad
       \Gamma,q(z)=0\vdash q(Sz)=0}
      {\Gamma\vdash q(t)=0}.
 \tag{ZI}\label{rule:ur-zero-ind}
\]
Only coverage, separation, selector computation, equality reasoning,
and UR were needed.  No predecessor, pairing, multiplication,
general primitive-recursion compiler, or faithful equality
characteristic was used.

The distinction between zero equations and arbitrary equations is
essential.  Coverage decides whether a single number is zero.  It
does not itself give a term $\delta(a,b)$ satisfying
$\delta(a,b)=0\leftrightarrow a=b$ for arbitrary $a,b$.

\subsection{The remaining bootstrap and the comparison boundary}

If the UR-based arithmetic could independently construct a scalar
term $\delta(a,b)$ and prove
\[
 \delta(a,b)=0\quad\longleftrightarrow\quad a=b
 \tag{EqRef}\label{eq:ur-missing-reflection}
\]
for actual equality throughout its domain, then ZI would recover
literal-equation induction.  Given an induction on $s(n)=t(n)$,
apply ZI to $\delta(s(n),t(n))=0$ and transport its base, step,
and conclusion through \eqref{eq:ur-missing-reflection}.
Together with the earlier proof that equation induction implies UR,
this would establish equivalence of the two arithmetic principles
over the stated common background.

The existing scalar proof of \eqref{eq:ur-missing-reflection}
cannot simply be imported at this point.  Its predecessor and
reflection developments use literal-equation inductions with
nonconstant right sides, including reconstruction invariants.
Using those results to justify a replacement for that very induction
would be circular unless their proofs were translated to UR or ZI.
The argument above removes part of the obstacle by deriving ZI, but
does not complete that translation.

There is a separate logical bookkeeping issue for a proposed native
$Q,D$ kernel.  The proofs in this section use ordinary intuitionistic
equality logic, contextual case distinction, substitution, and cut.
When replacing a primitive rule in the native kernel, the required
native versions of these resources must be derived without appealing
to the removed rule.  The ordinary arithmetic lemma alone is not a
certificate of equivalence between the complete native calculi.

Finally, a premise with no local induction hypothesis need not produce
a smaller implementation.  UR replaces the hypothesis extension by
construction of $f(t)$, $f(0)$, and an iterator with update $h$.
Its complete syntax checks, eigenvariable checks, logical rules, and
public theorem interface would all have to be measured.  No such
complete smaller checker or completed bi-interpretation is asserted
here.  The precise further research problem is to finish the
equality-reflection bootstrap under UR, then perform a matched
implementation comparison of the resulting complete calculus.


\section{Further research questions}

The following questions are deliberately separated from established claims.
\begin{enumerate}
\item Can equality reflection for arbitrary scalar equations be proved
from scalar iteration uniqueness, coverage, separation, and the two
computation laws, without invoking the equation-induction rule being
replaced? The retraction theorem reduces the final induction step to
that concrete bootstrap problem.
\item Can the predecessor development be completed under uniqueness with
a fully explicit list of common recurrences, then extended through the
residual equality-reflection invariant? This requires checking each
dependency rather than importing a classical equivalence theorem for a
different primitive-recursive language.
\item If the uniqueness route succeeds mathematically, does its complete
standalone checker improve the current size coordinates? Removing an
induction hypothesis alone does not answer this: its conclusion constructs
several term instances and an iterator.
\item Can a restriction of numerical cases avoid the product countermodel
while still deriving all native formula cases, with less checking
machinery? A successful proposal must explain where global logical case
distinction enters; equational predecessor identities alone do not do so.
\item Can the full typed arity checker be connected to the existing Lean
native semantics by a proved decoder and certificate-checking theorem?
That would strengthen the verification status of the whole family beyond
the written preservation proofs and regression checks supplied here.
\end{enumerate}

\section{Deliverable inventory and precise limits}

The archive contains this compiled report and its complete LaTeX source,
the three standalone checker implementations, exact matched controls,
inherited test certificates, the preservation proof, reproducible metrics
and tests, deterministic gzip-9 artifacts, selected Lean source snapshots,
the inherited scalar-to-iPRA written proofs, a provenance ledger, and a
SHA-256 manifest.

The new verified result is the 881/862-byte implementation reduction with
the stated matching contracts. The mathematical rule inventory in
Section~\ref{sec:presentation} is inherited. The product obstruction is a
proved rejection of a proposed weakening; the uniqueness results are
proved conditional progress toward another presentation. Neither is
represented as a completed simpler iPRA axiomatization. Absolute
minimality and a new formal proof of Python correctness remain unclaimed.

\begin{thebibliography}{9}
\bibitem{tait} W. W. Tait, \emph{Primitive Recursive Arithmetic and its
Role in the Foundations of Arithmetic: Historical and Philosophical
Reflections}, author manuscript, Section 2 (pp. 2--3).
\url{https://home.uchicago.edu/~wwtx/PRA2.pdf}.
\bibitem{readme} Vladimir Reshetnikov, \emph{Concise presentations of
intuitionistic primitive recursive arithmetic}, ProveIt README, pinned
commit \texttt{56249a4263cae9aad9ba87d23d0ee6d4ff6483bc}.
\url{https://github.com/VladimirReshetnikov/ProveIt/blob/56249a4263cae9aad9ba87d23d0ee6d4ff6483bc/Logic/PrimitiveRecursiveArithmetic/ConciseIPRA/README.md}.
\bibitem{native} ProveIt, \emph{Native.lean}, same snapshot.
\url{https://github.com/VladimirReshetnikov/ProveIt/blob/56249a4263cae9aad9ba87d23d0ee6d4ff6483bc/Logic/PrimitiveRecursiveArithmetic/ConciseIPRA/Lean/ConciseIPRA/Native.lean}.
\bibitem{inherited} ProveIt, \emph{NativeInherited.lean}, same snapshot;
theorem \texttt{iprf\_iff\_nprf}. See also the README's
\texttt{zprf\_iff\_nprf} comparison and the written source-restriction
proofs in the inherited packages.
\url{https://github.com/VladimirReshetnikov/ProveIt/blob/56249a4263cae9aad9ba87d23d0ee6d4ff6483bc/Logic/PrimitiveRecursiveArithmetic/ConciseIPRA/Lean/ConciseIPRA/NativeInherited.lean}.
\bibitem{bridge} ProveIt, \emph{NativeBiInterpretation.lean} and
\emph{BiInterpretation.lean}, same snapshot.
\url{https://github.com/VladimirReshetnikov/ProveIt/blob/56249a4263cae9aad9ba87d23d0ee6d4ff6483bc/Logic/PrimitiveRecursiveArithmetic/ConciseIPRA/Lean/ConciseIPRA/NativeBiInterpretation.lean}.
\bibitem{validated} ProveIt incoming archive
\texttt{ipra\_validated\_core\_2026-11-11.zip}, Git blob
\texttt{63b42195dc702cf115fbb9f472e96355427c9e95}.
\url{https://github.com/VladimirReshetnikov/ProveIt/blob/56249a4263cae9aad9ba87d23d0ee6d4ff6483bc/Logic/PrimitiveRecursiveArithmetic/ConciseIPRA/docs/incoming/ipra_validated_core_2026-11-11.zip}.
\bibitem{checked} ProveIt incoming archive
\texttt{ipra\_checked\_refinement\_2026-11-11.zip}, Git blob
\texttt{0583b5629b493c316d01b85d410961bf3e32866b}.
\url{https://github.com/VladimirReshetnikov/ProveIt/blob/56249a4263cae9aad9ba87d23d0ee6d4ff6483bc/Logic/PrimitiveRecursiveArithmetic/ConciseIPRA/docs/incoming/ipra_checked_refinement_2026-11-11.zip}.
\bibitem{onesided} ProveIt incoming archive
\texttt{ipra\_one\_sided\_transport\_2026-11-11.zip}, Git blob
\texttt{b3633919ebdc8ba0f66f4f522af9b595dc583bff}.
\url{https://github.com/VladimirReshetnikov/ProveIt/blob/56249a4263cae9aad9ba87d23d0ee6d4ff6483bc/Logic/PrimitiveRecursiveArithmetic/ConciseIPRA/docs/incoming/ipra_one_sided_transport_2026-11-11.zip}.
\end{thebibliography}
\end{document}
